\section{How it solves}
\label{sec:solver}

CFOP is four stages.  Each one is a separate module, each returns structured data,
and the pipeline verifies the whole solution before handing it back.

\subsection{The cross}

The four bottom edges are solved \emph{optimally}.  Only three things matter about
a cross: which four of the twelve edge positions the bottom edges occupy, which of
the four pieces is in which position, and how each is oriented.  That gives
\[
  \binom{12}{4} \cdot 4! \cdot 2^4 = 495 \cdot 24 \cdot 16 = 190{,}080
\]
states.  A breadth-first search over that space yields a distance table of one byte
per state, about a megabyte, and a solution is read off by walking downhill one move
at a time.  The table takes under a second to build, once per process.

No cross needs more than eight moves, and because the table stores exact distances,
the reported cross is as short as it possibly can be.  \texttt{--cross best} simply
repeats the search for all six colors and keeps the shortest.

\subsection{The first two layers}

The four corner--edge pairs are inserted one at a time, in the order front-right,
front-left, back-left, back-right, each insertion keeping the cross and every pair
already inserted.

Two mechanisms cooperate:

\begin{description}[leftmargin=0em,style=nextline]
  \item[The 41 case table] \texttt{data/f2l.json} holds the standard cases, each
  expressed in the front-right frame.  A slot is recognized by the position and
  orientation of its two pieces, and an algorithm is applied --- conjugated by a
  whole-cube turn when the slot is not the front-right one, so the display stays
  cross-down.
  \item[The verified search] Anything the table does not cover --- which includes
  any configuration where one of the pieces is stuck in another slot --- is inserted
  by an iterative deepening search over the cube, guided by three admissible
  heuristics: the distance to insert the pair, the distance to restore the cross and
  the distance to restore each pair already inserted.  The search only stops at a
  state where all of them are solved, so it cannot break an earlier stage.
\end{description}

A table result is always checked against the same goal before it is accepted; if it
would disturb something, the search takes over.  The pair subproblem itself is
small --- eight corner positions times three orientations times twelve edge
positions times two orientations, 576 states --- which is what makes the heuristic
cheap to evaluate at every node.

\subsection{Orientation and permutation of the last layer}

OLL and PLL are lookup tables of 57 and 21 cases.  Recognition does not use
hand-written pattern descriptions, which are easy to get subtly wrong.  Instead the
pattern of a case is \emph{computed from the case}: for each algorithm the engine
applies its inverse to a solved cube, reads off the pattern of the resulting state,
and stores that as the recognition key.  Table and recognizer therefore cannot
drift apart.

A case is only defined up to an adjustment of the top face, so every case is stored
for all sixteen combinations of a turn \emph{before} the algorithm and a turn
\emph{after} it --- the two-sided AUF of the cubing literature.  That is what makes
the tables complete.  The test suite proves it by enumerating the whole last layer:

\begin{center}
\begin{tabular}{@{}lll@{}}
\toprule
Space & Size & Verified \\
\midrule
Orientation states (OLL) & $3^3 \cdot 2^3 = 216$ & every state recognized and oriented \\
Permutation states (PLL) & $4! \cdot 4! / 2 = 288$ & every state recognized and permuted \\
\bottomrule
\end{tabular}
\end{center}

The 216 and 288 include the already-solved configuration, which needs no algorithm:
the pipeline reports \texttt{OLL (skip)} or \texttt{PLL (skip)} instead.

\subsection{Scrambles: a two-phase solver}

Random-state scrambles need a solver that can handle any state quickly, so the
project includes a Kociemba-style \emph{two-phase} solver.  Six coordinates describe
the cube compactly:

\begin{center}
\begin{tabular}{@{}llr@{}}
\toprule
Coordinate & Meaning & Values \\
\midrule
twist        & orientation of the eight corners & $3^7 = 2187$ \\
flip         & orientation of the twelve edges & $2^{11} = 2048$ \\
slice        & which positions hold the four slice edges & $\binom{12}{4} = 495$ \\
corner perm  & permutation of the corners & $8! = 40{,}320$ \\
edge perm    & permutation of the eight up/down edges & $8! = 40{,}320$ \\
slice perm   & permutation of the four slice edges & $4! = 24$ \\
\bottomrule
\end{tabular}
\end{center}

Phase~1 brings twist, flip and slice home using all eighteen face turns; phase~2
solves the rest using only half turns of R, L, F, B and any turn of U and D, which
cannot disturb what phase~1 fixed.  Both phases are iterative deepening searches
guided by breadth-first pruning tables (the worst phase-1 distances are 9, 9, 14
and 12).  A solve takes about 20 to 22 moves and, once the tables are built, a
fraction of a second.

To produce a scramble the generator draws a uniformly random solvable state, solves
it, and returns the \emph{inverse} of that solution --- so applying the scramble to a
solved cube really does produce the drawn state.  States that would need fewer than
two moves are redrawn, as Regulation 4b3 requires.

\subsection{Every solution is verified}

The pipeline concatenates the stages and applies the result to the scramble.  If the
cube is not solved, it raises an error rather than returning the solution: a
solution that does not check out is a bug, not a result.  The same flag is exposed as
\texttt{verified} in the JSON payload.

\subsection{Measured performance}

Figures from the development machine, for a random-state scramble of 20 to 22 moves:

\begin{center}
\begin{tabular}{@{}lr@{}}
\toprule
Step & Typical cost \\
\midrule
Build the scrambler tables (first use only) & 10 s \\
Generate a random-state scramble & 0.1--0.5 s \\
Cross, all six colors & under 1 ms \\
Four F2L pairs & 0.5--4 s (search-bound; see below) \\
OLL and PLL & under 5 ms \\
Whole solve & 0.7--4 s, 54 to 62 moves \\
\bottomrule
\end{tabular}
\end{center}

The F2L search dominates.  It is only used for configurations the 41-case table does
not cover, and it is correct for every one of them; a complete pair table would make
it unnecessary.  This is the single largest performance improvement available and is
listed again in Section~\ref{sec:trouble}.
